% ============================================================
\section{Discussão}\label{sec:discussao}
% ============================================================

Esta seção será escrita depois dos resultados e da robustez. Ela não traz números novos: junta
o que as Seções~\ref{sec:resultados} e~\ref{sec:robustez} mostram e responde às perguntas.

\previsto{
\begin{itemize}[leftmargin=1.2em]
    \item \textbf{Resposta a cada pergunta (P1--P4)}, confrontada com a hipótese
    correspondente (H1--H4) e com o grau de confiança dado pela robustez.
    \item \textbf{O que as redes mostram além das estatísticas de pares.} Por exemplo, se a
    queda da autossimilaridade~\cite{ethayarajh2019} vem acompanhada de troca de vizinhos
    locais, de novos \emph{hubs} ou de comunidades organizadas por contexto.
    \item \textbf{Papel da fragmentação do Qwen3.} O que muda entre palavras inteiras e
    fragmentos, e o que isso implica para a generalização a outros tokenizers.
    \item \textbf{Bloco 36 e previsão do próximo token.} Se a última camada se organiza pelo
    próximo token, e como isso se relaciona com os \emph{embeddings} amarrados
    (Seção~\ref{sec:mod-escolha}) e com a extensão ``vizinhos no vocabulário''.
    \item \textbf{Rede do vocabulário.} O que a geometria do espaço de entrada inteiro diz sobre
    os tipos que aparecem no português e sobre tokens pouco treinados.
    \item \textbf{Comparação com o piloto} (GPT-2 small, 104 tokens): o que se confirmou em
    escala e o que era artefato do desempate por posição.
\end{itemize}}
